\documentclass[10pt,a4paper]{amsart}
\usepackage{amsmath,amssymb,amsthm,graphicx}
\usepackage{verbatim}
\newenvironment{reference}{\comment}{\endcomment}
%\newenvironment{reference}{}{}
\newenvironment{slogan}{\comment}{\endcomment}
\newenvironment{history}{\comment}{\endcomment}
% For commutative diagrams we use Xy-pic
\usepackage[all]{xy}

% We use 2cell for 2-commutative diagrams.
\xyoption{2cell}
% Theorem environments.
%
\theoremstyle{plain}
\newtheorem{theorem}[subsection]{Theorem}
\newtheorem{proposition}[subsection]{Proposition}
\newtheorem{lemma}[subsection]{Lemma}

\theoremstyle{definition}
\newtheorem{definition}[subsection]{Definition}
\newtheorem{example}[subsection]{Example}
\newtheorem{exercise}[subsection]{Exercise}
\newtheorem{situation}[subsection]{Situation}

\theoremstyle{remark}
\newtheorem{remark}[subsection]{Remark}
\newtheorem{remarks}[subsection]{Remarks}

\numberwithin{equation}{subsection}

% Macros
%
\def\lim{\mathop{\mathrm{lim}}\nolimits}
\def\colim{\mathop{\mathrm{colim}}\nolimits}
\def\Spec{\mathop{\mathrm{Spec}}}
\def\Hom{\mathop{\mathrm{Hom}}\nolimits}
\def\Ext{\mathop{\mathrm{Ext}}\nolimits}
\def\SheafHom{\mathop{\mathcal{H}\!\mathit{om}}\nolimits}
\def\SheafExt{\mathop{\mathcal{E}\!\mathit{xt}}\nolimits}
\def\Sch{\mathit{Sch}}
\def\Mor{\mathop{\mathrm{Mor}}\nolimits}
\def\Ob{\mathop{\mathrm{Ob}}\nolimits}
\def\Sh{\mathop{\mathit{Sh}}\nolimits}
\def\NL{\mathop{N\!L}\nolimits}
\def\CH{\mathop{\mathrm{CH}}\nolimits}
\def\proetale{{pro\text{-}\acute{e}tale}}
\def\etale{{\acute{e}tale}}
\def\QCoh{\mathit{QCoh}}
\def\Ker{\mathop{\mathrm{Ker}}}
\def\Im{\mathop{\mathrm{Im}}}
\def\Coker{\mathop{\mathrm{Coker}}}
\def\Coim{\mathop{\mathrm{Coim}}}

% Boxtimes
%
\DeclareMathSymbol{\boxtimes}{\mathbin}{AMSa}{"02}

%
% Macros for moduli stacks/spaces
%
\def\QCohstack{\mathcal{QC}\!\mathit{oh}}
\def\Cohstack{\mathcal{C}\!\mathit{oh}}
\def\Spacesstack{\mathcal{S}\!\mathit{paces}}
\def\Quotfunctor{\mathrm{Quot}}
\def\Hilbfunctor{\mathrm{Hilb}}
\def\Curvesstack{\mathcal{C}\!\mathit{urves}}
\def\Polarizedstack{\mathcal{P}\!\mathit{olarized}}
\def\Complexesstack{\mathcal{C}\!\mathit{omplexes}}
% \Pic is the operator that assigns to X its picard group, usage \Pic(X)
% \Picardstack_{X/B} denotes the Picard stack of X over B
% \Picardfunctor_{X/B} denotes the Picard functor of X over B
\def\Pic{\mathop{\mathrm{Pic}}\nolimits}
\def\Picardstack{\mathcal{P}\!\mathit{ic}}
\def\Picardfunctor{\mathrm{Pic}}
\def\Deformationcategory{\mathcal{D}\!\mathit{ef}}
\usepackage{cleveref}
\crefname{originaldefinition}{original definition}{original definitions}
\Crefname{originaldefinition}{Original definition}{Original definitions}
\makeatletter
\newlabel{definition-fill-in-diagram}{{definition-fill-in-diagram}{source}{}{sourcepointer-1}{}}
\newlabel{definition-fill-in-diagram@cref}{{[originaldefinition][0][]definition-fill-in-diagram}{[1][0][]source}{}{}{}}
\makeatother
\crefname{originalequation}{original equation}{original equations}
\Crefname{originalequation}{Original equation}{Original equations}
\makeatletter
\newlabel{equation-diagram}{{equation-diagram}{source}{}{sourcepointer-2}{}}
\newlabel{equation-diagram@cref}{{[originalequation][0][]equation-diagram}{[1][0][]source}{}{}{}}
\makeatother
\crefname{originallemma}{original lemma}{original lemmas}
\Crefname{originallemma}{Original lemma}{Original lemmas}
\makeatletter
\newlabel{lemma-base-change-existence}{{lemma-base-change-existence}{source}{}{sourcepointer-3}{}}
\newlabel{lemma-base-change-existence@cref}{{[originallemma][0][]lemma-base-change-existence}{[1][0][]source}{}{}{}}
\makeatother
\crefname{originallemma}{original lemma}{original lemmas}
\Crefname{originallemma}{Original lemma}{Original lemmas}
\makeatletter
\newlabel{lemma-cat-dotted-arrows}{{lemma-cat-dotted-arrows}{source}{}{sourcepointer-4}{}}
\newlabel{lemma-cat-dotted-arrows@cref}{{[originallemma][0][]lemma-cat-dotted-arrows}{[1][0][]source}{}{}{}}
\makeatother
\crefname{originallemma}{original lemma}{original lemmas}
\Crefname{originallemma}{Original lemma}{Original lemmas}
\makeatletter
\newlabel{lemma-diagonal-diagonal}{{lemma-diagonal-diagonal}{source}{}{sourcepointer-5}{}}
\newlabel{lemma-diagonal-diagonal@cref}{{[originallemma][0][]lemma-diagonal-diagonal}{[1][0][]source}{}{}{}}
\makeatother
\crefname{originallemma}{original lemma}{original lemmas}
\Crefname{originallemma}{Original lemma}{Original lemmas}
\makeatletter
\newlabel{lemma-first-diagonal-separated-second-diagonal-closed}{{lemma-first-diagonal-separated-second-diagonal-closed}{source}{}{sourcepointer-6}{}}
\newlabel{lemma-first-diagonal-separated-second-diagonal-closed@cref}{{[originallemma][0][]lemma-first-diagonal-separated-second-diagonal-closed}{[1][0][]source}{}{}{}}
\makeatother
\crefname{originalreference}{original reference}{original references}
\Crefname{originalreference}{Original reference}{Original references}
\makeatletter
\newlabel{spaces-morphisms-lemma-quasi-compact-existence-universally-closed}{{spaces-morphisms-lemma-quasi-compact-existence-universally-closed}{source}{}{sourcepointer-7}{}}
\newlabel{spaces-morphisms-lemma-quasi-compact-existence-universally-closed@cref}{{[originalreference][0][]spaces-morphisms-lemma-quasi-compact-existence-universally-closed}{[1][0][]source}{}{}{}}
\makeatother
\crefname{originalreference}{original reference}{original references}
\Crefname{originalreference}{Original reference}{Original references}
\makeatletter
\newlabel{stacks-properties-section-properties-morphisms}{{stacks-properties-section-properties-morphisms}{source}{}{sourcepointer-8}{}}
\newlabel{stacks-properties-section-properties-morphisms@cref}{{[originalreference][0][]stacks-properties-section-properties-morphisms}{[1][0][]source}{}{}{}}
\makeatother
\title{If the categories of dotted arrows are setoids then the diagonal is separated}
\author{The Stacks Project Authors}
\date{stacks-project:stacks-morphisms.tex:local-snapshot-20261012; GNU FDL 1.2 or later}
\begin{document}
\maketitle
\noindent\textit{Editorial note.} All delivered mathematical text is the author's own, copied from the pinned Stacks Project chapter source sources/stacks-algebraic-stacks-morphisms/pinned-original-source.tex; no statement or proof step was rewritten, reordered or removed.\par\medskip
\noindent\textit{Editorial note.} Preamble provenance: the Stacks Project keeps one single shared preamble for every chapter; this card takes it directly from the staged sources/stacks-algebraic-stacks-morphisms/preamble.tex as three verbatim spans (lines 8-12, 14-18 and 154-217), so no derived file was written and nothing was added or redefined. The wrapper class and the base packages are supplied by the assembler.\par\medskip
\noindent\textit{Editorial note.} Bibliography provenance: neither the counted statement nor its proof contains a citation (this lane rescans the whole span before writing the card), so the card carries no bibliography and no entry was reconstructed.\par\medskip
\noindent\textit{Editorial note.} Selection provenance: this card was opened under the file-based control contract (CONTROL.md), with the mechanical pre-card gate run on its statement and proof interval before the card was written.\par\medskip
\noindent\textit{Allowed standard prerequisites (editorial).}\begin{itemize}
\item The second diagonal of a morphism of algebraic stacks through its inertia, and the criteria expressing separatedness of the diagonal through the closedness of that diagonal. Invoked at the start of the counted proof (source lines 9470-9477) as pointers to lemmas of this chapter. Reduces the assertion to the universal closedness of the second diagonal, which is quasi-compact by hypothesis.
\item The valuative criterion for the second diagonal, involving a 2-commutative diagram from a fraction field and a choice of 2-morphism, and the description of the categories of dotted arrows. The statement refers to the relevant definition and diagram, and the proof uses the corresponding lemma of this chapter (source lines 9480-9511). Supplies the shape of the diagrams to be filled and the automorphisms to be compared.
\end{itemize}\medskip
\ifdefined\hypertarget\hypertarget{sourcepointer-1}{}\fi
\noindent\textit{Source reference.} Original definition definition-fill-in-diagram: sources/stacks-algebraic-stacks-morphisms/pinned-original-source.tex, source lines 9116--9116. Navigation pointer only: the counted text refers to the Stacks item definition-fill-in-diagram, which belongs to another chapter of the project and is not part of this card. The pointer carries the target's own label as its number because this chapter's staged snapshot does not contain the other chapter's tag table; its content is not reproduced here.\par\medskip
\ifdefined\hypertarget\hypertarget{sourcepointer-2}{}\fi
\noindent\textit{Source reference.} Original equation equation-diagram: sources/stacks-algebraic-stacks-morphisms/pinned-original-source.tex, source lines 9080--9080. Navigation pointer only: the counted text refers to the Stacks item equation-diagram, which belongs to another chapter of the project and is not part of this card. The pointer carries the target's own label as its number because this chapter's staged snapshot does not contain the other chapter's tag table; its content is not reproduced here.\par\medskip
\ifdefined\hypertarget\hypertarget{sourcepointer-3}{}\fi
\noindent\textit{Source reference.} Original lemma lemma-base-change-existence: sources/stacks-algebraic-stacks-morphisms/pinned-original-source.tex, source lines 9514--9514. Navigation pointer only: the counted text refers to the Stacks item lemma-base-change-existence, which belongs to another chapter of the project and is not part of this card. The pointer carries the target's own label as its number because this chapter's staged snapshot does not contain the other chapter's tag table; its content is not reproduced here.\par\medskip
\ifdefined\hypertarget\hypertarget{sourcepointer-4}{}\fi
\noindent\textit{Source reference.} Original lemma lemma-cat-dotted-arrows: sources/stacks-algebraic-stacks-morphisms/pinned-original-source.tex, source lines 9447--9447. Navigation pointer only: the counted text refers to the Stacks item lemma-cat-dotted-arrows, which belongs to another chapter of the project and is not part of this card. The pointer carries the target's own label as its number because this chapter's staged snapshot does not contain the other chapter's tag table; its content is not reproduced here.\par\medskip
\ifdefined\hypertarget\hypertarget{sourcepointer-5}{}\fi
\noindent\textit{Source reference.} Original lemma lemma-diagonal-diagonal: sources/stacks-algebraic-stacks-morphisms/pinned-original-source.tex, source lines 3959--3959. Navigation pointer only: the counted text refers to the Stacks item lemma-diagonal-diagonal, which belongs to another chapter of the project and is not part of this card. The pointer carries the target's own label as its number because this chapter's staged snapshot does not contain the other chapter's tag table; its content is not reproduced here.\par\medskip
\ifdefined\hypertarget\hypertarget{sourcepointer-6}{}\fi
\noindent\textit{Source reference.} Original lemma lemma-first-diagonal-separated-second-diagonal-closed: sources/stacks-algebraic-stacks-morphisms/pinned-original-source.tex, source lines 1816--1816. Navigation pointer only: the counted text refers to the Stacks item lemma-first-diagonal-separated-second-diagonal-closed, which belongs to another chapter of the project and is not part of this card. The pointer carries the target's own label as its number because this chapter's staged snapshot does not contain the other chapter's tag table; its content is not reproduced here.\par\medskip
\ifdefined\hypertarget\hypertarget{sourcepointer-7}{}\fi
\noindent\textit{Source reference.} Original reference spaces-morphisms-lemma-quasi-compact-existence-universally-closed: sources/stacks-algebraic-stacks-morphisms/pinned-original-source.tex, source lines 9527--9527. Navigation pointer only: the counted text refers to the Stacks item spaces-morphisms-lemma-quasi-compact-existence-universally-closed, which belongs to another chapter of the project and is not part of this card. The pointer carries the target's own label as its number because this chapter's staged snapshot does not contain the other chapter's tag table; its content is not reproduced here.\par\medskip
\ifdefined\hypertarget\hypertarget{sourcepointer-8}{}\fi
\noindent\textit{Source reference.} Original reference stacks-properties-section-properties-morphisms: sources/stacks-algebraic-stacks-morphisms/pinned-original-source.tex, source lines 248--248. Navigation pointer only: the counted text refers to the Stacks item stacks-properties-section-properties-morphisms, which belongs to another chapter of the project and is not part of this card. The pointer carries the target's own label as its number because this chapter's staged snapshot does not contain the other chapter's tag table; its content is not reproduced here.\par\medskip
\begin{lemma}
\label{lemma-setoids-and-diagonal}
Let $f : \mathcal{X} \to \mathcal{Y}$ be a morphism of algebraic stacks.
If $\Delta_f$ is quasi-separated and if for every diagram
(\ref{equation-diagram}) and choice of $\gamma$ as in
Definition \ref{definition-fill-in-diagram}
the category of dotted arrows
is a setoid, then $\Delta_f$ is separated.
\end{lemma}

\begin{proof}
We are going to write out a detailed proof, but we strongly urge the
reader to find their own proof, inspired by reading the argument
given in the proof of Lemma \ref{lemma-cat-dotted-arrows}.

\medskip\noindent
Assume $\Delta_f$ is quasi-separated and for every diagram
(\ref{equation-diagram}) and choice of $\gamma$ as in
Definition \ref{definition-fill-in-diagram}
the category of dotted arrows is a setoid.
By Lemma \ref{lemma-diagonal-diagonal} it suffices to show that
$e : \mathcal{X} \to \mathcal{I}_{\mathcal{X}/\mathcal{Y}}$
is a closed immersion. By
Lemma \ref{lemma-first-diagonal-separated-second-diagonal-closed}
it in fact suffices to show that $e = \Delta_{f, 2}$ is
universally closed.
Either of these lemmas tells us that $e = \Delta_{f, 2}$ is quasi-compact
by our assumption that $\Delta_f$ is quasi-separated.

\medskip\noindent
In this paragraph we will show that $e$ satisfies the existence
part of the valuative criterion. Consider a $2$-commutative solid diagram
$$
\xymatrix{
\Spec(K) \ar[r]_x \ar[d]_j & \mathcal{X} \ar[d]^e \\
\Spec(A) \ar[r]^{(a, \theta)} & \mathcal{I}_{\mathcal{X}/\mathcal{Y}}
}
$$
and let $\alpha : (a, \theta) \circ j \to e \circ x$ be any $2$-morphism
witnessing the $2$-commutativity of the diagram (we use $\alpha$ instead
of the letter $\gamma$ used in Definition \ref{definition-fill-in-diagram}).
Note that $f \circ \theta = \text{id}$; we will use this below.
Observe that $e \circ x = (x, \text{id}_x)$ and
$(a, \theta) \circ j = (a \circ j, \theta \star \text{id}_j)$.
Thus we see that $\alpha$ is a $2$-arrow $\alpha : a \circ j \to x$
compatible with $\theta \star \text{id}_j$ and $\text{id}_x$.
Set $y = f \circ x$ and $\beta = \text{id}_{f \circ a}$.
Reading the arguments given in the proof of
Lemma \ref{lemma-cat-dotted-arrows}
backwards, we see that $\theta$ is an automorphism of the
dotted arrow $(a, \alpha, \beta)$ with
$$
\gamma : y \circ j \to f \circ x
\quad\text{equal to}\quad
\text{id}_f \star \alpha : f \circ a \circ j \to f \circ x
$$
On the other hand, $\text{id}_a$ is an automorphism too, hence
we conclude $\theta = \text{id}_a$ from the assumption on $f$.
Then we can take as dotted arrow for the displayed diagram above
the morphism $a : \Spec(A) \to \mathcal{X}$ with $2$-morphisms
$(a, \text{id}_a) \circ j \to (x, \text{id}_x)$ given by $\alpha$
and $(a, \theta) \to e \circ a$ given by $\text{id}_a$.

\medskip\noindent
By Lemma \ref{lemma-base-change-existence} any base change of $e$
satisfies the existence part of the valuative criterion.
Since $e$ is representable by algebraic spaces, it suffices to
show that $e$ is universally closed after a base change
by a morphism $I \to \mathcal{I}_{\mathcal{X}/\mathcal{Y}}$
which is surjective and smooth and with $I$ an algebraic space
(see Properties of Stacks, Section
\ref{stacks-properties-section-properties-morphisms}).
This base change $e' : X' \to I'$ is a quasi-compact
morphism of algebraic spaces which
satisfies the existence part of the valuative criterion
and hence is universally closed by
Morphisms of Spaces, Lemma
\ref{spaces-morphisms-lemma-quasi-compact-existence-universally-closed}.
\end{proof}

\end{document}
